\section*{Project contract}
Implement \texttt{blit\_sprite} to copy a rectangular sprite into a framebuffer. Pass the framebuffer descriptor in \texttt{a0}, the sprite descriptor in \texttt{a1}, and signed 32-bit placement coordinates in \texttt{a2} and \texttt{a3}, sign extended to the full register width. The sprite's local cell (0,0) is placed at (x,y). Clip against the framebuffer's logical width and height. Coordinates grow rightward and downward.

Return in \texttt{v0} the number of visible, opaque source pixels selected for copying. Return in \texttt{v1} the number of visible source pixels, including transparent ones. An opaque source pixel counts even when the destination already contains its color. A fully invisible or empty rectangle returns two zeros.

Each descriptor starts with an eight-byte address field pointing to four-byte-aligned pixel storage, followed by three four-byte words: width at byte 8, height at byte 12, and stride in words at byte 16. Descriptors themselves are eight-byte aligned. The sprite descriptor has an additional transparency key at byte 20. Width and height are each zero through eight; stride is at least width and at most sixteen. Every logical pixel cell fits inside its supplied 128-word allocation, possibly beginning at a nonzero word offset. Source, destination, and descriptors occupy disjoint storage.

Transparency compares the complete 32-bit source word with the key. Copy opaque words unchanged. Do not write transparent destination cells, clipped cells, row padding, guards, descriptors, or source pixels, even temporarily. Source reads may cover any logical source pixels, including clipped-away ones, but not padding or guards. Destination reads may cover logical framebuffer pixels. A multiword access is allowed only if all its bytes satisfy these rules. Reads of the selected descriptor fields are allowed; other caller data is not scratch storage.

Preserve the full values of \texttt{s0} through \texttt{s7}, \texttt{fp}, and \texttt{sp}, return through the supplied \texttt{ra}, and produce no console output. Private stack storage is permitted. The caller supplies valid descriptors; this exercise does not require malformed-descriptor handling or overlapping-image support. Equivalent implementations are accepted without prescribed scratch registers or instruction counts.
